\begin{theorem}[Rado embedding]
Let $r$ be a quasi-order on $Q$.  If $r$ is well-quasi-ordered and the
first-level domination relation $\mathsf{SetDomination}(r)$ on
$\mathcal P(Q)$ is not well-quasi-ordered, then there is a map

\[
e:\mathsf{IncreasingPair}\longrightarrow Q
\]

such that $e$ is a relation embedding from the Rado order to $r$:

\[
\mathsf{RelationEmbedding}(\mathsf{RadoOrder},r,e).
\]
Equivalently, $e$ is injective and, for all increasing pairs $p,q$,
\[
\mathsf{RadoOrder}(p,q)\quad\Longleftrightarrow\quad r(e(p),e(q)).
\]
\end{theorem}
\begin{proof}
By the bad-witness characterization \noderef{PowerWqoBadPairSequence},
the hypothesis that $\mathsf{SetDomination}(r)$ is not well-quasi-ordered
gives a map $f:\mathsf{IncreasingPair}\to Q$ satisfying
\noderef{BadPairSequence}:
\[
 m<n<l\quad\Longrightarrow\quad
 \neg r\bigl(f(m,n),f(n,l)\bigr).
\]
We use the two finite homogeneous consequences of Nash--Williams,
namely \noderef{TripleHomogeneous} and
\noderef{QuadrupleHomogeneous}, exactly as in the paper.

Color every three-element subset $\{i,j,k\}$, written with
$i<j<k$, by $1$ when
\[
r\bigl(f(i,j),f(i,k)\bigr)
\]
holds, and by $0$ otherwise.  Apply \noderef{TripleHomogeneous} to the
whole set of natural numbers.  We obtain an infinite set $N$ on which
the color is constant.  The constant color cannot be $0$: fix any
$i\in N$ and enumerate the infinite tail $N/i$ increasingly, as allowed
by \noderef{InfiniteSetEnumeration}.  The resulting sequence
$t\mapsto f(i,n_t)$ has no forward related pair, because its two terms
at indices $u<v$ occur in the colored triple
$i<n_u<n_v$.  This contradicts $h_Q$, since
$\mathsf{WellQuasiOrdered}(r)$ means that every sequence has indices
$u<v$ with related values.  Thus every triple in $N$ has color $1$:
\[
i<j<k\text{ in }N\quad\Longrightarrow\quad
r\bigl(f(i,j),f(i,k)\bigr).
\tag{1}
\]

Now color every four-element subset $\{i,j,k,l\}$ of $N$, with
$i<j<k<l$, by $1$ when
\[
r\bigl(f(i,j),f(k,l)\bigr)
\]
holds, and by $0$ otherwise.  Applying
\noderef{QuadrupleHomogeneous} gives an infinite $M\subseteq N$ on
which this color is constant.  The constant cannot be $0$.  Indeed,
choose successively pairs
\[
i_0<j_0<i_1<j_1<i_2<j_2<\cdots
\]
from $M$.  If the color were $0$, then for every $u<v$ the quadruple
$i_u<j_u<i_v<j_v$ would give
\[
\neg r\bigl(f(i_u,j_u),f(i_v,j_v)\bigr),
\]
so the sequence $u\mapsto f(i_u,j_u)$ would contradict $h_Q$.  Hence
\[
i<j<k<l\text{ in }M\quad\Longrightarrow\quad
r\bigl(f(i,j),f(k,l)\bigr).
\tag{2}
\]

Let $a=\min M$ and put $X=M\setminus\{a\}$.  This is infinite; choose
an increasing enumeration $x:\mathbb N\to X$ using
\noderef{InfiniteSetEnumeration}.  For an increasing pair
$p=(i,j)$, define
\[
e(p)=f\bigl(x(i),x(j)\bigr).
\]
The inequalities needed to form the pair on the right follow from the
strict increase of $x$.

We first prove preservation of the Rado order, using the definition
\noderef{RadoOrder}.  If $p=(i,j)$ and $q=(k,l)$ have $i=k$ and
$j\leq l$, then either $j=l$, in which case reflexivity of the
quasi-order gives $r(e(p),e(q))$, or $j<l$, in which case (1) gives
$r(f(x(i),x(j)),f(x(i),x(l)))$.  If instead $j<k$, then
$x(i)<x(j)<x(k)<x(l)$, so (2) gives
$r(e(p),e(q))$.  Therefore
\[
\mathsf{RadoOrder}(p,q)\Longrightarrow r(e(p),e(q)).
\tag{3}
\]

For reflection, assume $r(e(p),e(q))$, where $p=(i,j)$ and $q=(k,l)$.
Suppose that $\mathsf{RadoOrder}(p,q)$ fails.  By the two clauses in
\noderef{RadoOrder}, this means
\[
k\leq j\quad\text{and}\quad (i\neq k\ \text{or}\ l<j).
\tag{4}
\]
There are three cases.  If $l<j$, choose any $n>j$.  The Rado order has
$(k,l)\leq(j,n)$ by its second clause.  Preservation (3) and
transitivity therefore give
\[
r\bigl(f(x(i),x(j)),f(x(j),x(n))\bigr),
\]
contradicting the badness of $f$ at the increasing triple
$x(i)<x(j)<x(n)$.

Assume next that $l\geq j$ and $i<k$.  Then
$(i,k)\leq(i,j)$ by the first Rado clause.  Preservation followed by
the assumed relation gives
\[
r\bigl(f(x(i),x(k)),f(x(k),x(l))\bigr),
\]
again contradicting badness at $x(i)<x(k)<x(l)$.

The remaining possibility in (4) is $l\geq j$ and $k<i$.  Since
$a<x(k)<x(i)$, the pair $(a,x(k))$ is an increasing pair and
$(a,x(k))\leq(x(i),x(j))$ by the second Rado clause.  Preservation and
transitivity now give
\[
r\bigl(f(a,x(k)),f(x(k),x(l))\bigr),
\]
contradicting badness at $a<x(k)<x(l)$.  All cases are impossible, so
\[
r(e(p),e(q))\Longrightarrow\mathsf{RadoOrder}(p,q).
\tag{5}
\]

It remains to check injectivity.  If $e(p)=e(q)$, reflexivity of $r$
gives both $r(e(p),e(q))$ and $r(e(q),e(p))$.  By (5), both Rado
relations hold.  The two-clause definition of the Rado order is
antisymmetric: a cross-clause relation $p\leq q$ cannot be reversed,
and if the first clause holds in both directions then the first
coordinates agree and antisymmetry of $\leq$ on $\mathbb N$ gives equality
of the second coordinates.  Hence $p=q$.  Thus $e$ is injective.

Combining injectivity with (3) and (5), and unfolding
\noderef{RelationEmbedding}, proves
\[
\mathsf{RelationEmbedding}(\mathsf{RadoOrder},r,e).
\]
This is the embedding asserted in lines 985--987 of the paper.
\end{proof}
